\clearpage\thispagestyle{empty}
\pdfbookmark[0]{Quellen zu Pallas (12. Dezember 1814)}{pallas1814dec-sources}
\begin{center}\textbf{Quellen zu Pallas (12. Dezember 1814)}\end{center}
\noindent Diese KI-Übertragung ins heutige Deutsch folgt dem Bericht mit der Überschrift Pallas in den \href{https://gdz.sub.uni-goettingen.de/id/PPN236007467}{\textit{Werken}, Band~VI} von Gauß (Göttingen, 1874), S.~379--381. Er erschien unter der zweiten Überschrift Göttingen in den \textit{Göttingischen gelehrten Anzeigen}, Stück~199 (12. Dezember 1814), S.~\href{https://images.p.sub.uni-goettingen.de/iiif/image/gdz:PPN319721507_1814:00001993/full/1600,/0/default.jpg}{1985}, \href{https://images.p.sub.uni-goettingen.de/iiif/image/gdz:PPN319721507_1814:00001994/full/1600,/0/default.jpg}{1986}, \href{https://images.p.sub.uni-goettingen.de/iiif/image/gdz:PPN319721507_1814:00001995/full/1600,/0/default.jpg}{1987}, \href{https://images.p.sub.uni-goettingen.de/iiif/image/gdz:PPN319721507_1814:00001996/full/1600,/0/default.jpg}{1988}, \href{https://images.p.sub.uni-goettingen.de/iiif/image/gdz:PPN319721507_1814:00001997/full/1600,/0/default.jpg}{1989}, \href{https://images.p.sub.uni-goettingen.de/iiif/image/gdz:PPN319721507_1814:00001998/full/1600,/0/default.jpg}{1990}. Die Links führen zum digitalisierten Exemplar der Niedersächsischen Staats- und Universitätsbibliothek Göttingen.
\par\smallskip\noindent In der Werkausgabe steht im Tabellenkopf Seberg. In der letzten Deklinationszelle zum 3. November steht 10 6 mit Zwischenraum, und die Differenz der mittleren Länge zur zweiten Opposition lautet $-12.9$. Im Erstdruck stehen an diesen Stellen Seeberg, 10,6 und $-2{,}9$. Der Dezemberbericht gibt in der Vorausbestimmung der Opposition 50.5 Sekunden an, der gesonderte Maibrief 50 Sekunden.
\par\smallskip\noindent \textit{Fußnote des Erstdrucks, S.~1988.} Durch einen Druckfehler wurde in Stück~129, S.~1282, ihre Polhöhe mit $54^\circ\,17'\,50''3$ statt mit $54^\circ\,42'\,50''3$ angegeben. Die Anmerkung bezieht sich auf die neue Königsberger Sternwarte und stammt aus dem Erstdruck. KI-Übertragung und KI-Satz.
\clearpage
